\documentclass[11pt]{article}

\usepackage[margin=1in]{geometry}
\usepackage{amsmath}
\usepackage{booktabs}
\usepackage{array}
\usepackage{enumitem}
\usepackage{graphicx}
\usepackage[hidelinks]{hyperref}
\usepackage[numbers,sort&compress]{natbib}

\title{Calibrated, Leakage-Safe Ranking of Short-Horizon Equity Setups:\\
An Honest Walk-Forward Study on NSE Stocks}

\author{
Kartik Gupta\thanks{Student author. Experiments, implementation, and first manuscript draft.}\\
\small [Department / Institution to be added]\\
\small [City, Country]\\
\and
Faculty Mentor Name\thanks{Faculty mentor/advisor. To be confirmed after review.}\\
\small [Department / Institution to be added]\\
\small [City, Country]
}

\date{\today}

\begin{document}
\maketitle

\begin{abstract}
Retail equity-screening systems often report attractive backtests while hiding
lookahead leakage, weak calibration, or overfit selection rules. This paper
presents a modest empirical study of a reproducible pipeline for ranking
short-horizon equity setups on Indian National Stock Exchange (NSE) stocks. For
each stock and trading day, the system constructs causal technical, relative
strength, market-breadth, sector, regime, and earnings-proximity features using
only information available up to the close of day $t$. A concrete trade is then
defined using an ATR-based entry band, target, stop, and approximately five-day
holding horizon. The label is produced by a triple-barrier rule: target first is
a win, stop first is a loss, and unresolved trades exit at the horizon. We train
a gradient-boosted tree classifier to estimate the probability that a setup hits
target before stop, calibrate this probability on a time-ordered holdout tail,
and evaluate all predictions using walk-forward testing with an embargo equal
to the holding period. The contribution is not a claim of a profitable trading
system. Instead, we ask whether calibrated, leakage-safe ranking improves over
the unranked base rate and rule-based ranking when measured honestly out of
sample. All performance numbers in this draft are placeholders until the
experiment scripts are run and independently checked.
\end{abstract}

\noindent\textbf{Keywords:} financial machine learning; probability calibration;
walk-forward validation; triple-barrier labeling; NSE equities; leakage-safe
backtesting.\\
\textbf{Suggested arXiv categories:} cs.LG, q-fin.ST, cs.CE.

\section{Introduction}

Short-horizon stock-selection tools are attractive because they convert noisy
market data into an apparently simple ranked list of candidates. The same
simplicity also makes them easy to overstate. A trading model can look useful if
it is evaluated on future information, tuned repeatedly on the same period,
or judged only by raw accuracy rather than by calibrated probabilities and
realistic trade outcomes. These failure modes are especially important for
student or retail research, where the goal should be reproducibility and honest
measurement rather than promotional claims.

This paper studies a single question:

\begin{quote}
\emph{Does calibrated, leakage-safe, walk-forward ranking of NSE equity setups
beat the unranked base rate and a rule-based ranking baseline out of sample,
and by how much when the trades are graded in R-multiples after costs?}
\end{quote}

The system under study is a daily pipeline that defines a concrete short-horizon
trade for each stock-day and estimates the probability that the trade reaches
its target before its stop. The model is not evaluated by whether it can predict
the next close exactly. It is evaluated by whether the ranked setups produce
better out-of-sample barrier outcomes than simpler baselines.

The intended contribution is deliberately modest:

\begin{enumerate}[leftmargin=*]
    \item We describe a leakage-safe data pipeline that creates causal
    stock-day features and triple-barrier labels for short-horizon NSE equity
    setups.
    \item We estimate and calibrate setup win probabilities using a
    time-ordered training and calibration design rather than random
    cross-validation.
    \item We report an honest walk-forward evaluation, including classification
    metrics, probability reliability, economic R-multiple metrics, baselines,
    ablations, sensitivity checks, and limitations.
\end{enumerate}

This work should not be read as financial advice or as evidence that the system
is safe for live trading. A realistic outcome for this type of model may be only
a thin edge, no edge, or an edge concentrated in specific market regimes. That
is precisely why the evaluation protocol is the main object of the paper.

\section{Related Work}

Machine learning has been applied widely to equity prediction and asset pricing.
Large-scale studies such as \citet{gu2020empirical} show that non-linear models
can extract predictive structure from financial characteristics, but also that
evaluation design matters. Technical-indicator based equity prediction is a
smaller and noisier setting: many candidate signals are correlated, non-stationary,
and sensitive to transaction costs.

This study follows several ideas from financial machine learning practice.
\citet{lopezdeprado2018afml} emphasizes meta-labeling, triple-barrier labeling,
and purged or embargoed validation to avoid overlap between training labels and
test periods. Backtest overfitting is a known risk in strategy research
\citep{bailey2014probability,bailey2017deflated}; therefore, this paper treats
parameter sweeps as sensitivity analysis rather than proof.

Probability calibration is also central. A model can rank examples reasonably
while producing probabilities that do not match realized frequencies.
\citet{platt1999probabilistic} introduced a widely used sigmoid calibration
approach for support vector machines, while \citet{niculescu2005predicting}
compare calibration methods such as Platt scaling and isotonic regression. We
use isotonic calibration because the desired output is an interpretable estimate
of setup win probability, not only a ranking score. Calibration is evaluated
using the Brier score \citep{brier1950verification} and reliability tables.
Discrimination is summarized using ROC-AUC \citep{fawcett2006roc}.

The model family is gradient-boosted decision trees, implemented using the
histogram-based gradient boosting classifier provided by scikit-learn
\citep{pedregosa2011scikit}. Gradient boosting is a strong baseline for tabular
data with non-linear interactions; related systems include LightGBM
\citep{ke2017lightgbm}. The contribution here is not a new model architecture.
It is the construction, calibration, and evaluation of a transparent empirical
pipeline for a specific stock-selection task.

\section{Data}

\subsection{Source and Universe}

The dataset is built from daily open, high, low, close, and volume (OHLCV) data
for Indian NSE equities, downloaded or cached through \texttt{yfinance}
\citep{aroussi2024yfinance}. The market context uses the Nifty 50 index
(\texttt{\^{}NSEI}). The current project uses approximately 474 NSE stocks for
training and a smaller large-cap daily-signal universe of approximately 80
stocks for deployment-style signal generation.

The approximate training table contains 491,000 stock-day rows over about 1,125
trading days. The unranked base win rate is approximately 31.6\%. Exact values
should be regenerated from the experiment scripts and inserted before
submission:
\texttt{[[RESULT: exact dataset rows]]},
\texttt{[[RESULT: exact symbol count]]},
\texttt{[[RESULT: exact date range]]}, and
\texttt{[[RESULT: exact base win rate]]}.

\subsection{Stock-Day Row Construction}

For every stock $i$ and trading day $t$, the pipeline creates one candidate row
using only information available up to the close of day $t$. Features are
computed from historical price, volume, index, sector, breadth, and earnings
calendar data. The label is computed from future bars $t+1,\ldots,t+H$, where
$H$ is the holding horizon. Features and labels are therefore separated in time:
past data defines the input, and future data defines the outcome.

\subsection{Preprocessing}

Rows with insufficient rolling-window history are skipped. Indicators requiring
20-day, 50-day, or 14-day windows are computed only after enough observations
exist. The index series is aligned by trading date. The system uses conservative
handling for ambiguous daily bars because daily OHLCV data cannot tell whether
the target or stop was hit first within the same session.

\section{Methodology}

\subsection{Feature Set}

Table~\ref{tab:features} summarizes the causal feature groups. Each feature is
computed as-of day $t$ only.

\begin{table}[ht]
\centering
\small
\caption{Causal feature groups used for each stock-day row.}
\label{tab:features}
\begin{tabular}{p{0.27\linewidth} p{0.63\linewidth}}
\toprule
\textbf{Group} & \textbf{Features and definition} \\
\midrule
Trend and momentum &
Close relative to EMA20; EMA20 relative to EMA50; 5-day and 20-day returns;
RSI(14); ADX(14). \\
\addlinespace
Volatility and range &
ATR as a percentage of close; daily high-low range as a percentage of close. \\
\addlinespace
Volume and liquidity &
Current volume divided by 20-day average volume; liquidity score used in final
ranking. \\
\addlinespace
Support and resistance &
Distance from 20-day resistance in ATR units; distance from 20-day support in
ATR units. \\
\addlinespace
Relative and sector strength &
Relative strength versus the market index; mean sector relative strength;
sector breadth relative to market breadth. \\
\addlinespace
Market regime and breadth &
Regime score; buy/sell regime alignment; fraction of universe above EMA50 and
EMA20; fraction advancing over five days; change in breadth. \\
\addlinespace
Rule-engine context &
Direction indicator; buy score; sell score; chart setup quality. \\
\addlinespace
Earnings proximity &
Days since earnings; clipped days until earnings; pre-earnings flag;
post-earnings flag. Future earnings dates are clipped to reduce leakage from
unscheduled future events. \\
\bottomrule
\end{tabular}
\end{table}

\subsection{Trade Definition and Triple-Barrier Label}

For each row, the rule engine defines a direction, entry band, target, and stop
using the latest close and average true range (ATR). For a buy setup, the target
is above the entry and the stop is below it; for a sell setup, the directions
are reversed. Let $R$ denote the risk distance from entry to stop. Outcomes are
expressed in R-multiples so trades with different prices and volatility can be
compared.

The label is assigned over the next $H$ trading sessions:
\begin{itemize}[leftmargin=*]
    \item \textbf{WIN} if the target is hit before the stop; the binary label is
    $y=1$ and realized R is the target-to-risk multiple.
    \item \textbf{LOSS} if the stop is hit before the target; the binary label is
    $y=0$ and realized R is $-1$.
    \item \textbf{TIME\_EXIT} if neither barrier is touched by the horizon; the
    binary label is conservatively treated as $y=0$, while realized R is marked
    to the horizon close.
\end{itemize}

If both target and stop are touched within the same daily bar, the default
offline experiment scores the trade as a loss. This is pessimistic, but it
avoids flattering the strategy when intraday ordering is unknown.

\subsection{Model}

The predictive model is a scikit-learn
\texttt{HistGradientBoostingClassifier}. It is chosen because gradient-boosted
trees are strong, CPU-efficient models for tabular data with non-linear
interactions and missingness-resistant splits. The model estimates
\[
    \hat{p}_{i,t} = P(y_{i,t}=1 \mid x_{i,t}),
\]
where $x_{i,t}$ is the causal feature vector for stock $i$ on day $t$.

The raw classifier output is calibrated with isotonic regression using
\texttt{CalibratedClassifierCV}. Calibration is fit on the time-ordered tail of
the training period, not on randomly shuffled folds. This preserves the temporal
structure of the problem.

\subsection{Expected R and Final Ranking}

For a candidate setup, expected R is computed as
\[
    E[R] =
    p_{\mathrm{win}} \cdot R_{\mathrm{target}}
    -
    p_{\mathrm{loss}} \cdot R_{\mathrm{stop}}
    -
    c_{\mathrm{txn}}
    -
    c_{\mathrm{slip}},
\]
where the default transaction and slippage costs are approximately
$0.03R$ and $0.02R$, respectively.

The deployed ranking score is a weighted combination of calibrated win
probability, expected R, historical setup win rate, regime alignment, chart
quality, relative strength, liquidity, model stability, and risk penalties:
\[
\begin{aligned}
\mathrm{final\_score} ={}&
0.30p_{\mathrm{win}}
+ 0.20E[R]
+ 0.15\mathrm{AdjSetupWinRate} \\
&+ 0.10\mathrm{RegimeAlignment}
+ 0.10\mathrm{ChartQuality}
+ 0.05\mathrm{RelativeStrength} \\
&+ 0.05\mathrm{Liquidity}
+ 0.05\mathrm{ModelStability}
- \mathrm{RiskPenalties}.
\end{aligned}
\]

For the balanced profile, a signal is shown only if it clears gates such as
calibrated $p_{\mathrm{win}} \ge 0.64$, risk--reward $\ge 1.55$, positive
expected R, and final score $\ge 0.58$. These deployment gates are evaluated
separately from the broader ranking study. Diversification caps highly
correlated selections, with pairwise recent-return correlation above 0.92
treated as a duplicate exposure.

\subsection{Walk-Forward Validation With Embargo}

All reported metrics are based on concatenated out-of-sample predictions from a
walk-forward design. The default protocol is:
\begin{enumerate}[leftmargin=*]
    \item Sort all rows by trading date.
    \item Use an initial warm-up period of approximately 500 trading days.
    \item Train only on rows strictly before the test window.
    \item Remove an embargo gap equal to the holding horizon ($H=5$ days) so
    training labels cannot overlap the test window.
    \item Fit calibration on the last 15\% of the training dates.
    \item Predict the next out-of-sample window, approximately 60 trading days.
    \item Slide forward and repeat until all available out-of-sample windows are
    predicted.
\end{enumerate}

This design is stricter than random cross-validation because no future row is
allowed to train a past prediction.

\section{Experiments and Results}

\subsection{Experiment Commands}

The following commands regenerate the main experiment outputs from the project
root after backend environment variables and caches are configured:

\begin{verbatim}
cd backend
python scripts/train_ml.py --hold 5 --top-n 5 --cost 0.05
python scripts/train_ml.py --hold 5 --top-n 10 --cost 0.05
python scripts/analyze_ml_oos.py
python scripts/run_backtest.py --risk Balanced --top-n 5 --cost 0.05 --hold 5
python scripts/run_backtest_sweep.py
\end{verbatim}

If caches need to be rebuilt, run:
\begin{verbatim}
cd backend
python scripts/train_ml.py --rebuild --hold 5 --top-n 5 --cost 0.05
\end{verbatim}

\subsection{Main Classification Metrics}

Table~\ref{tab:classification} reports out-of-sample classification quality.
These values must be filled from \path{scripts/train_ml.py}.

\begin{table}[ht]
\centering
\caption{Out-of-sample classification metrics.}
\label{tab:classification}
\begin{tabular}{lc}
\toprule
\textbf{Metric} & \textbf{Value} \\
\midrule
OOS rows & \texttt{[[RESULT: OOS row count]]} \\
OOS trading days & \texttt{[[RESULT: OOS trading days]]} \\
ROC-AUC & \texttt{[[RESULT: OOS ROC-AUC]]} \\
Brier score & \texttt{[[RESULT: OOS Brier score]]} \\
Base win rate, all rows & \texttt{[[RESULT: OOS base win rate]]} \\
\bottomrule
\end{tabular}
\end{table}

\subsection{Calibration Reliability}

Table~\ref{tab:reliability} compares predicted probability buckets with
realized win rates. A useful probability model should be monotonic and close to
the diagonal, although financial data may remain noisy even after calibration.

\begin{table}[ht]
\centering
\caption{Reliability table for calibrated probabilities.}
\label{tab:reliability}
\resizebox{\linewidth}{!}{%
\begin{tabular}{lccc}
\toprule
\textbf{Predicted probability bucket} & \textbf{Count} & \textbf{Mean predicted} & \textbf{Actual win rate} \\
\midrule
\texttt{[[RESULT: bucket 1 range]]} & \texttt{[[RESULT: bucket 1 n]]} & \texttt{[[RESULT: bucket 1 pred]]} & \texttt{[[RESULT: bucket 1 actual]]} \\
\texttt{[[RESULT: bucket 2 range]]} & \texttt{[[RESULT: bucket 2 n]]} & \texttt{[[RESULT: bucket 2 pred]]} & \texttt{[[RESULT: bucket 2 actual]]} \\
\texttt{[[RESULT: bucket 3 range]]} & \texttt{[[RESULT: bucket 3 n]]} & \texttt{[[RESULT: bucket 3 pred]]} & \texttt{[[RESULT: bucket 3 actual]]} \\
\texttt{[[RESULT: bucket 4 range]]} & \texttt{[[RESULT: bucket 4 n]]} & \texttt{[[RESULT: bucket 4 pred]]} & \texttt{[[RESULT: bucket 4 actual]]} \\
\texttt{[[RESULT: bucket 5 range]]} & \texttt{[[RESULT: bucket 5 n]]} & \texttt{[[RESULT: bucket 5 pred]]} & \texttt{[[RESULT: bucket 5 actual]]} \\
\bottomrule
\end{tabular}
}
\end{table}

\subsection{Economic Test}

The economic evaluation selects top-$N$ setups per day and grades realized R
after costs. We report decisive R separately from time exits because unresolved
trades marked to the horizon close can otherwise hide weak barrier performance.

\begin{table}[ht]
\centering
\caption{Baseline comparison for top-$N$ selection after costs.}
\label{tab:economic}
\small
\resizebox{\linewidth}{!}{%
\begin{tabular}{lrrrrrr}
\toprule
\textbf{Selection} & \textbf{Trades} & \textbf{Win/Loss/Timeout} & \textbf{Decisive R} & \textbf{Avg R} & \textbf{PF} & \textbf{Wilson LB} \\
\midrule
All signals, no selection & \texttt{[[RESULT: all trades]]} & \texttt{[[RESULT: all W/L/T]]} & \texttt{[[RESULT: all decisive R]]} & \texttt{[[RESULT: all avg R]]} & \texttt{[[RESULT: all PF]]} & \texttt{[[RESULT: all Wilson LB]]} \\
Rule-ranked top-$N$ & \texttt{[[RESULT: rule trades]]} & \texttt{[[RESULT: rule W/L/T]]} & \texttt{[[RESULT: rule decisive R]]} & \texttt{[[RESULT: rule avg R]]} & \texttt{[[RESULT: rule PF]]} & \texttt{[[RESULT: rule Wilson LB]]} \\
ML-ranked top-$N$ & \texttt{[[RESULT: ML trades]]} & \texttt{[[RESULT: ML W/L/T]]} & \texttt{[[RESULT: ML decisive R]]} & \texttt{[[RESULT: ML avg R]]} & \texttt{[[RESULT: ML PF]]} & \texttt{[[RESULT: ML Wilson LB]]} \\
\bottomrule
\end{tabular}
}
\end{table}

The honest interpretation should be filled after running the experiments:
\texttt{[[RESULT: honest verdict on whether ML beats rule baseline]]}. If the
ML-ranked row does not clearly beat the rule-ranked row after costs, the paper
should state that result directly.

\subsection{Ablations}

Table~\ref{tab:ablations} lists required ablations. These are intended to test
whether the result depends on calibration, breadth, earnings, and relative or
sector strength features.

\begin{table}[ht]
\centering
\caption{Ablation study.}
\label{tab:ablations}
\small
\resizebox{\linewidth}{!}{%
\begin{tabular}{lcccc}
\toprule
\textbf{Variant} & \textbf{ROC-AUC} & \textbf{Brier} & \textbf{Top-$N$ avg R} & \textbf{Top-$N$ decisive R} \\
\midrule
Full calibrated model & \texttt{[[RESULT]]} & \texttt{[[RESULT]]} & \texttt{[[RESULT]]} & \texttt{[[RESULT]]} \\
Uncalibrated GBM & \texttt{[[RESULT]]} & \texttt{[[RESULT]]} & \texttt{[[RESULT]]} & \texttt{[[RESULT]]} \\
Drop market/sector breadth & \texttt{[[RESULT]]} & \texttt{[[RESULT]]} & \texttt{[[RESULT]]} & \texttt{[[RESULT]]} \\
Drop earnings proximity & \texttt{[[RESULT]]} & \texttt{[[RESULT]]} & \texttt{[[RESULT]]} & \texttt{[[RESULT]]} \\
Drop relative/sector strength & \texttt{[[RESULT]]} & \texttt{[[RESULT]]} & \texttt{[[RESULT]]} & \texttt{[[RESULT]]} \\
Logistic regression baseline & \texttt{[[RESULT]]} & \texttt{[[RESULT]]} & \texttt{[[RESULT]]} & \texttt{[[RESULT]]} \\
\bottomrule
\end{tabular}
}
\end{table}

\subsection{Sensitivity Checks}

The paper should include sensitivity to top-$N$, costs, and holding horizon.
Table~\ref{tab:sensitivity} is a placeholder for these results.

\begin{table}[ht]
\centering
\caption{Sensitivity checks.}
\label{tab:sensitivity}
\small
\resizebox{\linewidth}{!}{%
\begin{tabular}{lcccc}
\toprule
\textbf{Setting} & \textbf{Trades} & \textbf{Avg R} & \textbf{Decisive R} & \textbf{Profit factor} \\
\midrule
Top-3, cost 0.05R & \texttt{[[RESULT]]} & \texttt{[[RESULT]]} & \texttt{[[RESULT]]} & \texttt{[[RESULT]]} \\
Top-5, cost 0.05R & \texttt{[[RESULT]]} & \texttt{[[RESULT]]} & \texttt{[[RESULT]]} & \texttt{[[RESULT]]} \\
Top-10, cost 0.05R & \texttt{[[RESULT]]} & \texttt{[[RESULT]]} & \texttt{[[RESULT]]} & \texttt{[[RESULT]]} \\
Top-5, cost 0.03R & \texttt{[[RESULT]]} & \texttt{[[RESULT]]} & \texttt{[[RESULT]]} & \texttt{[[RESULT]]} \\
Top-5, cost 0.10R & \texttt{[[RESULT]]} & \texttt{[[RESULT]]} & \texttt{[[RESULT]]} & \texttt{[[RESULT]]} \\
Hold horizon 3 days & \texttt{[[RESULT]]} & \texttt{[[RESULT]]} & \texttt{[[RESULT]]} & \texttt{[[RESULT]]} \\
Hold horizon 10 days & \texttt{[[RESULT]]} & \texttt{[[RESULT]]} & \texttt{[[RESULT]]} & \texttt{[[RESULT]]} \\
\bottomrule
\end{tabular}
}
\end{table}

\subsection{Figures}

The final paper should include the following figures. The boxes below are
placeholders so the draft compiles before images are generated.

\begin{figure}[ht]
\centering
\fbox{\begin{minipage}[c][1.3in][c]{0.78\linewidth}
\centering
\texttt{[[FIGURE: calibration / reliability diagram]]}\\
Generate from \texttt{backend/scripts/train\_ml.py} OOS predictions.
\end{minipage}}
\caption{Predicted probability versus realized win rate.}
\label{fig:calibration}
\end{figure}

\begin{figure}[ht]
\centering
\fbox{\begin{minipage}[c][1.3in][c]{0.78\linewidth}
\centering
\texttt{[[FIGURE: distribution of realized R and cumulative R]]}\\
Generate from \texttt{backend/scripts/ml\_oos.pkl}.
\end{minipage}}
\caption{Realized R distribution and cumulative R for selected trades.}
\label{fig:r_distribution}
\end{figure}

\begin{figure}[ht]
\centering
\fbox{\begin{minipage}[c][1.3in][c]{0.78\linewidth}
\centering
\texttt{[[FIGURE: sensitivity curve for cost or top-N]]}\\
Generate from top-$N$ and cost sweeps.
\end{minipage}}
\caption{Sensitivity of selected-trade R metrics to selection depth or costs.}
\label{fig:sensitivity}
\end{figure}

\section{Discussion and Limitations}

Several limitations should be kept explicit.

\textbf{First, daily bars cannot resolve intraday ordering.} If target and stop
are both touched within one daily candle, the true first hit is unknown. The
offline experiment treats such cases conservatively, but this may still differ
from live trade execution.

\textbf{Second, the universe may contain survivorship and coverage bias.} The
training set uses symbols available in the cached data. If delisted or formerly
liquid stocks are missing, historical performance may be overstated.

\textbf{Third, the result may be regime dependent.} A model that performs only
in bullish periods is not necessarily a stock-picking model; it may be a market
regime filter. The analysis must therefore report slices by year, direction,
and regime.

\textbf{Fourth, transaction costs are simplified.} A flat R-based cost is useful
for comparison, but actual execution costs vary by spread, liquidity, slippage,
order type, and market conditions.

\textbf{Fifth, this is not a live trading study.} The paper evaluates historical
out-of-sample predictions produced by a walk-forward procedure. Live deployment
can introduce data delays, corporate-action issues, behavioral selection, and
execution errors that are not captured here.

\textbf{Finally, calibration can drift.} The feedback loop records realized
outcomes of live signals, but maintaining calibration over time requires regular
monitoring. A positive result in one historical sample should be treated as a
hypothesis for further validation, not as a permanent property of the market.

\section{Conclusion and Future Work}

This paper presents a reproducible, leakage-safe framework for evaluating
calibrated short-horizon equity setup ranking on NSE stocks. The central claim
is methodological: stock-selection systems should be judged with causal
features, future-only labels, time-ordered calibration, walk-forward testing,
embargoed splits, realistic costs, and transparent baselines. Whether the
current model adds a statistically and economically meaningful edge remains to
be filled from the experiment outputs. If the result is weak or regime-specific,
that is still a useful research outcome because it prevents overclaiming.

Future work includes adding intraday data to resolve target-stop ordering,
testing delisted-symbol coverage, expanding to other markets, comparing more
model classes, measuring calibration drift over live paper-trading periods, and
formalizing uncertainty intervals for top-$N$ economic metrics.

\section*{Acknowledgements}

I thank \texttt{[[MENTOR NAME]]} for reviewing the research question,
methodology, and manuscript. Any remaining mistakes are my own.

\section*{Data and Code Availability}

The code for the pipeline and experiments is intended to be released at
\url{https://github.com/Kartik962005/ai-trading-mvp}. Daily OHLCV data are
obtained through \texttt{yfinance}/Yahoo Finance subject to their terms of use.
Generated caches and experiment outputs should be documented with exact command
lines, random seeds, library versions, and train/test split dates before
submission.

\section*{Educational-Use Disclaimer}

This paper is for educational and research purposes only. It is not investment
advice, financial advice, or a recommendation to buy, sell, or short any
security. Historical backtests and walk-forward evaluations do not guarantee
future results.

\section*{Numbers to Fill From Experiments}

\begin{enumerate}[leftmargin=*]
    \item Dataset summary: rows, symbols, date range, trading days, exact base
    win rate. Command: \texttt{python scripts/train\_ml.py --hold 5 --top-n 5 --cost 0.05}.
    \item Main OOS metrics: ROC-AUC, Brier score, OOS row count, OOS trading
    days. Command: same as above.
    \item Reliability table: predicted-probability bucket ranges, counts, mean
    predicted probability, actual win rate. Command: same as above; copy
    ``Calibration reliability'' output.
    \item Economic comparison: ML top-$N$, rule top-$N$, and all-signals rows:
    trades, W/L/timeout counts, decisive net R, average net R, profit factor,
    Wilson lower bound, total R. Command: same as above, with \texttt{--top-n 5}
    and \texttt{--top-n 10}.
    \item Regime/year/direction slices. Command:
    run \texttt{analyze\_ml\_oos.py}.
    \item Rule-engine baseline report. Command:
    run \texttt{run\_backtest.py} with Balanced risk, top-5, cost 0.05R,
    and hold horizon 5.
    \item Sensitivity checks for top-$N$ and cost assumptions. Commands:
    rerun \texttt{train\_ml.py} with \texttt{--top-n 3}, \texttt{--top-n 5},
    \texttt{--top-n 10}, and costs \texttt{0.03}, \texttt{0.05}, and
    \texttt{0.10}.
    \item Ablation results: calibrated vs uncalibrated, drop breadth, drop
    earnings, drop relative/sector strength, logistic regression baseline.
    These may require adding ablation flags/scripts before submission.
    \item Figures: calibration diagram, realized R distribution/cumulative R,
    and top-$N$ or cost sensitivity curve.
    \item Final honest verdict: state whether ML-ranked selection beats the
    rule-ranked baseline out of sample after costs. Do not soften a negative or
    inconclusive result.
\end{enumerate}

\begin{thebibliography}{99}

\bibitem[Lopez de Prado(2018)]{lopezdeprado2018afml}
Lopez de Prado, M. (2018).
\newblock \emph{Advances in Financial Machine Learning}.
\newblock Wiley.

\bibitem[Gu et~al.(2020)]{gu2020empirical}
Gu, S., Kelly, B., and Xiu, D. (2020).
\newblock Empirical asset pricing via machine learning.
\newblock \emph{The Review of Financial Studies}, 33(5):2223--2273.

\bibitem[Platt(1999)]{platt1999probabilistic}
Platt, J.~C. (1999).
\newblock Probabilistic outputs for support vector machines and comparisons to
regularized likelihood methods.
\newblock In \emph{Advances in Large Margin Classifiers}, pages 61--74. MIT
Press.

\bibitem[Niculescu-Mizil and Caruana(2005)]{niculescu2005predicting}
Niculescu-Mizil, A. and Caruana, R. (2005).
\newblock Predicting good probabilities with supervised learning.
\newblock In \emph{Proceedings of the 22nd International Conference on Machine
Learning}, pages 625--632.

\bibitem[Brier(1950)]{brier1950verification}
Brier, G.~W. (1950).
\newblock Verification of forecasts expressed in terms of probability.
\newblock \emph{Monthly Weather Review}, 78(1):1--3.

\bibitem[Fawcett(2006)]{fawcett2006roc}
Fawcett, T. (2006).
\newblock An introduction to ROC analysis.
\newblock \emph{Pattern Recognition Letters}, 27(8):861--874.

\bibitem[Pedregosa et~al.(2011)]{pedregosa2011scikit}
Pedregosa, F., Varoquaux, G., Gramfort, A., Michel, V., Thirion, B.,
Grisel, O., Blondel, M., Prettenhofer, P., Weiss, R., Dubourg, V.,
Vanderplas, J., Passos, A., Cournapeau, D., Brucher, M., Perrot, M., and
Duchesnay, E. (2011).
\newblock Scikit-learn: Machine learning in Python.
\newblock \emph{Journal of Machine Learning Research}, 12:2825--2830.

\bibitem[Ke et~al.(2017)]{ke2017lightgbm}
Ke, G., Meng, Q., Finley, T., Wang, T., Chen, W., Ma, W., Ye, Q., and
Liu, T.-Y. (2017).
\newblock LightGBM: A highly efficient gradient boosting decision tree.
\newblock In \emph{Advances in Neural Information Processing Systems}, volume
30.

\bibitem[Wilson(1927)]{wilson1927probable}
Wilson, E.~B. (1927).
\newblock Probable inference, the law of succession, and statistical inference.
\newblock \emph{Journal of the American Statistical Association},
22(158):209--212.

\bibitem[Bailey et~al.(2017)]{bailey2014probability}
Bailey, D.~H., Borwein, J.~M., Lopez de Prado, M., and Zhu, Q.~J. (2017).
\newblock The probability of backtest overfitting.
\newblock \emph{Journal of Computational Finance}, 20(4):39--69.

\bibitem[Bailey and Lopez de Prado(2014)]{bailey2017deflated}
Bailey, D.~H. and Lopez de Prado, M. (2014).
\newblock The deflated Sharpe ratio: Correcting for selection bias, backtest
overfitting, and non-normality.
\newblock \emph{Journal of Portfolio Management}, 40(5):94--107.

\bibitem[Tashman(2000)]{tashman2000outofsample}
Tashman, L.~J. (2000).
\newblock Out-of-sample tests of forecasting accuracy: An analysis and review.
\newblock \emph{International Journal of Forecasting}, 16(4):437--450.

\bibitem[Aroussi(2024)]{aroussi2024yfinance}
Aroussi, R. (2024).
\newblock yfinance.
\newblock \url{https://github.com/ranaroussi/yfinance}.

\end{thebibliography}

\end{document}
